\section{Positioned homologies in characteristic at least five}\label{sec:higher-homologies}
This section locates the scalar elements of the cubic boundary in one arrangement fibre; their precise positions will allow the full higher-dimensional group to be generated.

Assume throughout Sections~\ref{sec:higher-homologies}--\ref{sec:higher-amalgamation} that $p\ge5$, $n\ge3$ and $d\ge n+2$, and put $m=d-3$. The exact sequences of Section~\ref{sec:cubic-cech} give an actual extension of finite local systems
\begin{equation}\label{eq:higher-extension}
0\longrightarrow P\longrightarrow V_{\mathrm{bdry}}
 \longrightarrow E=\bigoplus_{|I|=n-2}E_I\longrightarrow0.
\end{equation}
Here $E_I=H^1_{\mathrm{et}}(Q\cap\Lambda_I,\mathbf F_p)$. Since $3$ is invertible, each three-point permutation module splits canonically into its constant line and its sum-zero plane. Every incoming differential of the degree-zero intersection complex lands in the constant lines. Thus the second filtration sequence in Section~\ref{sec:cubic-cech-quotients} splits naturally as
\begin{equation}\label{eq:higher-P}
P=P_{\mathrm{const}}\oplus\bigoplus_{|J|=n-1}Q_J,
\qquad Q_J=\mathbf F_p^{Z_J}/\mathbf F_p\mathbf1.
\end{equation}
The constant part has dimension $\binom mn$. This is a splitting of the actual local system, not just a decomposition of its semisimplification.

These sequences extend to the parameter open allowing the cubic itself to be a transverse union of three hyperplanes. Indeed its positive-dimensional strata remain connected; cubic strata of dimension at least two have no positive coherent cohomology by the hypersurface sequence, while its plane sections are ordinary nodal triangles. Its zero-dimensional strata remain three distinct points. The same intersection-complex calculation therefore gives locally free cohomology of the same ranks and invertible Frobenius. Section~\ref{sec:cubic-generic} supplies a finite \'etale Lang system throughout this open. Restriction to the smooth-cubic open leaves its finite image unchanged, by the connected-cover argument there.

Let $H$ be this boundary image. Theorem~\ref{thm:joint} and \eqref{eq:higher-P} give
\begin{equation}\label{eq:higher-group-extension}
1\longrightarrow N\longrightarrow H\longrightarrow
 A=\prod_J S_3\times\prod_I\mathbf F_p^\times\longrightarrow1,
\qquad N\subset\operatorname{Hom}(E,P).
\end{equation}
The kernel acts identically on $P$ and $E$, so it is an additive $\mathbf F_p$-subspace of the displayed Hom space. The order of $A$ is prime to $p$.

There is also a distinguished subgroup $\Gamma\simeq S_3$ in $H$: at a cubic consisting of three general hyperplanes, permute those three unlabelled factors and leave all other factors labelled. The connected ordered configuration cover gives this action on the arrangement fibre. Its action on each $E_I$ is sign, since the cohomology of a triangle is its oriented cycle line. It acts on each $Q_J$ as the standard two-dimensional representation and fixes $P_{\mathrm{const}}$. In particular $P$ has no sign constituent for $\Gamma$.

\subsection{A complement containing the arrangement permutations}
The group extension \eqref{eq:higher-group-extension} has a complement containing $\Gamma$. Here is an elementary verification, including this last qualification. A set-theoretic section has a $2$-cocycle $c(a,b)\in N$. Set
\[
t(a)=|A|^{-1}\sum_{b\in A}c(a,b).
\]
Summing the cocycle identity in the third argument gives
\[
c(a,b)=a\,t(b)-t(ab)+t(a).
\]
Changing the section by $-t$ makes it a group homomorphism. The image of $\Gamma$ in $A$ is injective. The difference between its given lift and its lift in this complement is a $1$-cocycle on $\Gamma$ with coefficients in $N$. Averaging this cocycle expresses it as a coboundary, because $6$ is invertible. Conjugating the complement by the corresponding element of $N$ makes it contain $\Gamma$.

The scalar subgroup on each individual $E_I$ commutes with the image of $\Gamma$ in $A$. Its lift in this complement therefore commutes with $\Gamma$. The $\Gamma$-sign isotypic subspace of $V_{\mathrm{bdry}}$ maps isomorphically onto $E$ in \eqref{eq:higher-extension}, because $P$ has no sign constituent. The lifted scalar acts as the chosen scalar on that one line, as the identity on the other lines, and as the identity on $P$. This identifies its position, beyond the conjugacy class obtained in Proposition~\ref{prop:determinant-reflections}.

\subsection{Exterior coordinates at an arrangement}
Let $E_0=\mathbf F_p^d$ have basis $e_1,\ldots,e_d$, put $u=\sum_i e_i$, and set $W=E_0/\mathbf F_pu$. At a transverse arrangement of $d$ hyperplanes, the incidence resolution gives the $S_d$-equivariant marking
\begin{equation}\label{eq:higher-wedge}
V=\bigwedge^nW.
\end{equation}
Indeed its last differential is $u\wedge-:\bigwedge^{n-1}E_0\to\bigwedge^nE_0$, and there is no next term. The quotient is \eqref{eq:higher-wedge}, of dimension $\binom{d-1}n$. All cohomology is the stable constant incidence part, so this is an ordinary fibre of the same Lang system.

Fix a triple $T=\{a,b,c\}$ and define
\[
A_T=\langle e_a-e_c,e_b-e_c\rangle,
\qquad B_T=\langle e_i:i\notin T\rangle.
\]
Since $3$ is invertible, $W=A_T\oplus B_T$. Permutations of $T$ act standardly on $A_T$ and trivially on $B_T$. Let $S_T$ be their sign isotypic subspace in $V$, and let $P_T$ be the kernel of the sign projector. Then
\begin{equation}\label{eq:higher-block}
S_T=\bigwedge^2A_T\otimes\bigwedge^{n-2}B_T,
\qquad V=S_T\oplus P_T.
\end{equation}
For $I\subset\{1,\ldots,d\}\setminus T$ of size $n-2$, put
\[
v_{T,I}=(e_a-e_c)\wedge(e_b-e_c)\wedge e_I.
\]
These form a basis of $S_T$. Refining the cubic into its three components identifies their images in the elliptic quotient in \eqref{eq:higher-extension}: the cochain supported on $I$ and two of $a,b,c$ is the alternating triangle cochain. Its projection to $E_J$ is zero for $J\ne I$ and is three times an oriented triangle generator for $J=I$. One can read this directly in the double intersection complex: the degree-one refinement term over the plane indexed by $I$ is its triangle complex, and no other outside index set occurs in that term. Since $3$ is invertible, this gives exactly the indicated coordinate lines.

Let $\phi_{T,I}$ be the coordinate dual of $v_{T,I}$ on $S_T$, extended by zero on $P_T$. The scalar construction above supplies the exact elements
\begin{equation}\label{eq:higher-homology}
R_{T,I}(z)=1+(z-1)v_{T,I}\phi_{T,I},
\qquad z\in\mathbf F_p^\times,
\end{equation}
in the geometric image. Conjugation by the arrangement $S_d$ gives them for every triple $T$.

\section{Elementary matrices on the alternating blocks}\label{sec:higher-block-generation}
This section uses explicit overlap coefficients to turn the positioned homologies into a full linear group on each $S_T$.

Fix $T=\{a,b,c\}$, write $k=n-2$, and choose adjacent $k$-subsets of its complement,
\[
I=K\cup\{i\},\qquad L=K\cup\{j\},\qquad |K|=k-1.
\]
Use compatible orientations, placing $e_i,e_j$ before $e_K$. Set
\[
\begin{gathered}
v=v_{T,I},\quad v'=v_{T,L},\quad
T'=\{a,b,j\},\quad J=K\cup\{c\},\quad w=v_{T',J},\\
\phi=\phi_{T,I},\quad\phi'=\phi_{T,L},\quad\psi=\phi_{T',J}.
\end{gathered}
\]
The required coefficients are
\begin{equation}\label{eq:higher-overlap}
\phi(w)=-\frac13,\qquad\psi(v)=0,
\qquad\psi(v')=-1,\qquad\phi(v')=0.
\end{equation}
For a complete calculation, put $A_0=e_a-e_c$, $B_0=e_b-e_c$ and $s=\sum_{r\notin T}e_r$. Then $e_c=-(A_0+B_0+s)/3$. The projection onto $S_T$ of
\[
((e_a-e_j)\wedge(e_b-e_j))\wedge e_c\wedge e_K
\]
is
\[
-\frac13 A_0\wedge B_0\wedge s\wedge e_K
-\frac23 A_0\wedge B_0\wedge e_j\wedge e_K.
\]
Its coefficients at $v,v'$ are $-1/3,-1$. Interchanging $c,j$ in the same calculation gives $\psi(v')=-1$. Finally every row of the coordinate determinant defining $\psi$ vanishes on $e_i$, since $i\notin T'\cup J$, so $\psi(v)=0$. This proves \eqref{eq:higher-overlap}, with any additional outside indices present.

Choose $a_0,b_0\ne0,1$. Conjugate $R_w(b_0)$ by $R_v(a_0)$ and multiply by $R_w(b_0)^{-1}$. Since $\psi(v)=0$, the result is
\[
1+(1-b_0^{-1})(a_0-1)\phi(w)\,v\psi.
\]
Its coefficient is nonzero, and powers give $1+t v\psi$ for every $t\in\mathbf F_p$. Conjugating by $R_{v'}(c_0)$, with $c_0\ne0,1$, fixes the centre $v$ and changes its row by
\[
(c_0^{-1}-1)\psi(v')\phi'.
\]
Multiplication by the inverse original transvection gives a nonzero multiple of $v\phi'$. Powers again supply
\begin{equation}\label{eq:higher-elementary}
1+t v_{T,I}\phi_{T,L}\qquad(t\in\mathbf F_p).
\end{equation}
The graph of $k$-subsets joined by a one-index replacement is connected. Commutators of consecutive elementary matrices along a simple path give every ordered pair of distinct basis indices. Together with the diagonal scalings \eqref{eq:higher-homology}, they generate
\begin{equation}\label{eq:higher-full-block}
G_T=\operatorname{GL}(S_T)\quad\hbox{acting identically on }P_T.
\end{equation}
Here $1\le k\le m-1$, so every block has dimension at least two.

\section{Amalgamating the blocks}\label{sec:higher-amalgamation}
This section proves an elementary generation lemma and verifies its hypotheses for all higher-dimensional arrangement blocks, including when $p$ divides $d$.

\begin{lemma}\label{lem:higher-amalgamation}
Let $p$ be odd and let $V=S_i\oplus P_i$ be a collection of decompositions over $\mathbf F_p$, with $\dim S_i\ge2$. Put $L_i=\operatorname{Ann}(P_i)\subset V^*$. Suppose the graph joining $i,j$ when $S_i\cap S_j\ne0$ is connected, the $S_i$ span $V$, and the $L_i$ span $V^*$. Then the groups $\operatorname{GL}(S_i)$ acting identically on $P_i$ generate $\operatorname{GL}(V)$.
\end{lemma}
\begin{proof}
For $v\ne0$ and a row space $L$ evaluating nontrivially on $v$, define
\[
\mathcal A(v,L)=\{1+v\phi:\phi\in L,\ 1+\phi(v)\ne0\}.
\]
If $L,M$ both evaluate nontrivially on $v$, the two axis groups generate $\mathcal A(v,L+M)$. To prove this, choose $\lambda\in L$, $\mu\in M$ with $\lambda(v)=\mu(v)=1$. All unipotent rows in $L\cap\operatorname{Ann}(v)$ and $M\cap\operatorname{Ann}(v)$ are present. For $a\ne0,1$,
\[
[1+(a-1)v\lambda][1+(a-1)v\mu]^{-1}
 =1+(a-1)v(\lambda-\mu).
\]
Powers give every $\mathbf F_p$ multiple of this row difference. These three spaces span $(L+M)\cap\operatorname{Ann}(v)$, and the original axis scalars supply every nonzero value of $1+\phi(v)$. This proves the assertion.

Start with one block. Its group supplies $\mathcal A(v,L_i)$ for every nonzero $v\in S_i$. Inductively suppose this holds on every block of a connected collection with the full row space $L_C=\sum_{i\in C}L_i$. Add a block $S_j$ meeting $S_i$ in a nonzero $v$. The axis assertion first gives $\mathcal A(v,L_C+L_j)$. Each old block group preserves $L_C+L_j$: if $g$ acts identically on $P_i$, then $\phi g-\phi\in L_i$. Transitivity of $\operatorname{GL}(S_i)$ on its nonzero vectors supplies the enlarged axis group for every centre in $S_i$. Propagate along the old connected intersection graph, and use the new block group to obtain it on $S_j$ too. Induction gives axis rows from the total sum $V^*$ at every nonzero centre in every block.

Choose a basis of $V$ from the union of the $S_i$. All elementary transvections and all individual diagonal scalings in that basis are now present. They generate $\operatorname{GL}(V)$.
\end{proof}

We verify the three hypotheses for \eqref{eq:higher-full-block}. Intrinsically,
\[
S_T=(\bigwedge^2 A_T)\wedge\bigwedge^{n-2}W.
\]
If $T,T'$ share two vertices, $A_T+A_{T'}$ has dimension three. The nonzero space
\[
\bigwedge^3(A_T+A_{T'})\wedge\bigwedge^{n-3}W
\]
lies in $S_T\cap S_{T'}$. The graph of triples joined by a two-vertex overlap is connected. This proves the intersection hypothesis.

Let $D\subset W$ be spanned by all differences $e_i-e_j$. Its codimension is at most one. The spaces $\bigwedge^2 A_T$ span $\bigwedge^2D$, as is seen from the triples $\{i,j,d\}$. Hence the $S_T$ span
\[
(\bigwedge^2D)\wedge\bigwedge^{n-2}W.
\]
Every basis wedge of degree $n\ge3$ contains at least $n-1\ge2$ factors from a codimension-at-most-one subspace. The displayed span is therefore all $\bigwedge^nW$. This argument includes $p\mid d$, when $D$ has codimension one.

For the dual assertion, put $r=d-1$ and $q=r-n$. Exterior duality gives
\[
V^*\simeq\bigwedge^qW\otimes(\det W)^*.
\]
On the permutations of $T$, the determinant of $W$ is sign. Thus the sign isotypic space $L_T=\operatorname{Ann}(P_T)$ corresponds to the invariant subspace in $\bigwedge^qW$, tensored with $(\det W)^*$. Because $W=A_T\oplus B_T$ and $p>3$, that invariant subspace is exactly $\bigwedge^qB_T$: $A_T$ has no invariants and $\bigwedge^2A_T$ is sign. The spaces $\bigwedge^qB_T$ span $\bigwedge^qW$. Indeed $q=d-1-n\le d-4$, and every wedge of $q$ coordinate classes lies outside some triple $T$. Those coordinate wedges span the whole exterior power.

Lemma~\ref{lem:higher-amalgamation} now gives the full geometric $\operatorname{GL}_{\binom{d-1}n}(\mathbf F_p)$. The arithmetic image contains it and lies in the same group, so is equal. This proves Theorem~\ref{thm:higher-images} for $p\ge5$.
